% phdid: partial homogeneity in staggered difference-in-differences
% Draft article for the Stata Journal.
%
% Typeset with the Stata Journal class, which is distributed by the
% Journal and is not included here:
%     https://www.stata-journal.com/submissions/
% Place sj.cls and stata.sty beside this file, then
%     pdflatex phdid_sj
%
% Copyright (c) 2026 Rohan Wagle and Parush Arora.  MIT licensed.

\documentclass{sj}
\usepackage{amsmath,amssymb,booktabs,graphicx}

\stabstract{%
In a staggered difference-in-differences design the treatment effect is a
vector of cohort-time effects, not a single number.  Estimating every one
separately is unbiased but inefficient when some are in fact equal; pooling
them into a single two-way fixed-effects coefficient is efficient but biased
whenever the heterogeneity is genuine.  We describe \texttt{phdid}, which
treats the choice between these corners as a partition-selection problem on
the cohort-time cells and implements the two estimators of
\citet{AroraWagle2026}: a Dirichlet process mixture fitted by collapsed Gibbs
sampling, and an $\ell_0$-penalised estimator with a two-step agglomerative
search.  Because the estimators depend on the data only through the
first-stage coefficient vector and its covariance, \texttt{phdid} works after
any heterogeneity-robust first stage that posts \texttt{e(b)} and
\texttt{e(V)}.  The package also provides specification tests for whether a
partition is recoverable at all, exact partition enumeration for small
designs, regularisation paths, and sampler diagnostics.}

\title[phdid: Partial homogeneity in staggered DiD]{Partial homogeneity in
staggered difference-in-differences}

\author{Rohan Wagle\\Department of Economics\\Ashoka University\\Sonipat,
India\\rohan.wagle\_ug2023@ashoka.edu.in
\and Parush Arora\\Department of Economics\\Ashoka University\\Sonipat,
India\\parush.arora@ashoka.edu.in}

\keywords{\inserttag, difference-in-differences, staggered adoption,
treatment effect heterogeneity, Dirichlet process, partition selection, model
selection, Bayesian nonparametrics}

\begin{document}
\maketitle

\section{Introduction}

A staggered difference-in-differences design with $G$ treated cohorts and $T$
periods identifies not one treatment effect but $K$ of them, one per
cohort-time cell $(g,t)$ with $t \ge g$.  The modern estimators
\citep{CallawaySantAnna2021,Wooldridge2025} return all $K$, and correctly
warn against pooling them with a single two-way fixed-effects coefficient when
the effects differ.

That warning leaves an unanswered question.  The $K$ cell estimates are
individually noisy, and in applied work many of them are not meaningfully
different from one another.  Reporting all $K$ is unbiased but wasteful;
reporting one is efficient but wrong.  The useful object lies between: a
partition of the cells into $m < K$ groups within which the effects are
genuinely equal.  \citet{AroraWagle2026} call this \emph{partial homogeneity},
show that under the true partition the grouped estimator is unbiased and
strictly more efficient than the flexible one, and give two ways to recover
the partition when it is unknown.

This article describes \texttt{phdid}, a Stata implementation of both.  The
package is deliberately a second-stage tool: it does not estimate treatment
effects from micro data.  Section~\ref{sec:method} explains why that is not a
limitation but a theorem, and what follows from it for the design of the
command.

\section{Methods and formulas}
\label{sec:method}

\subsection{The model}

Let $\tau = (\tau_1,\dots,\tau_K)'$ collect the true cohort-time average
treatment effects and let $\hat\tau$ be any asymptotically normal first-stage
estimator with
\[
  \hat\tau \;\sim\; \mathcal{N}(\tau, \Sigma),
\]
$\Sigma$ the \emph{joint} covariance.  A partition $\mathcal{C} =
\{C_1,\dots,C_m\}$ of $\{1,\dots,K\}$ imposes $\tau_k = \phi_p$ for all $k \in
C_p$.  Writing $R$ for the $K \times m$ indicator matrix of $\mathcal{C}$ and
$\Omega = \Sigma^{-1}$, the restricted generalised least squares fit is
\begin{equation}
  \hat\phi = (R'\Omega R)^{-1} R'\Omega\hat\tau,
  \qquad
  \hat\tau_{\mathcal{C}} = R\hat\phi,
  \qquad
  \widehat{\mathrm{Var}}(\hat\tau_{\mathcal{C}}) = R (R'\Omega R)^{-1} R',
  \label{eq:gls}
\end{equation}
with deviance $D(\mathcal{C}) = \hat\tau'\Omega\hat\tau - \hat\phi'R'\Omega
\hat\tau$ against the flexible fit.

\subsection{Why the first stage is a sufficient statistic}

Everything in \eqref{eq:gls}, and everything in the marginal likelihood of the
Bayesian model below, enters only through $R'\Omega\hat\tau$ and $R'\Omega R$.
This is Lemma~1 of \citet{AroraWagle2026}.  Two first stages that produce the
same $(\hat\tau, \Sigma)$ therefore produce the same partition posterior,
whatever micro data lie behind them.

For a Stata package this is an unusually convenient theorem, because
$(\hat\tau, \Sigma)$ is exactly what every estimation command in Stata already
posts as \texttt{e(b)} and \texttt{e(V)}.  \texttt{phdid setup} reads those
and nothing else.  The package therefore composes with \texttt{csdid},
\texttt{jwdid}, \texttt{did\_imputation}, or a matrix the user assembled
themselves, with no reimplementation and no assumption about how the first
stage was computed.

\subsection{Remark 1: the covariance is not diagonal}

Cohort-time cells share units and share the unit and time fixed effects, so
$\Sigma$ has substantial off-diagonal mass.  Discarding it is not a harmless
speed-up.  In the design of \citet{AroraWagle2026} it understates the
posterior variance of aggregates by a factor of about 3.5, and moves
individual co-clustering probabilities by up to 0.39.  \texttt{phdid} uses the
exact $\Sigma$ throughout, warns when handed a diagonal one, and provides
\texttt{phdid covcheck} to quantify the difference on the user's own design.

\subsection{The $\ell_0$-penalised estimator}

The estimator of Appendix~B of \citet{AroraWagle2026} minimises
$D(\mathcal{C}) + \lambda \,\mathrm{pen}(\mathcal{C})$ in two steps.
Step~1 is a greedy agglomerative search from $K$ singletons, merging at each
stage the pair minimising
\begin{equation}
  \Delta(A,B) \;=\; \frac{n_A n_B}{n_A + n_B}\,(\bar\tau_A - \bar\tau_B)^2,
  \label{eq:merge}
\end{equation}
with $n$ a precision weight.  Step~2 re-estimates by \eqref{eq:gls} under the
selected partition.  Step~2 is not cosmetic: the criterion \eqref{eq:merge}
treats the design as orthogonal, and only the re-estimated fit is best linear
unbiased when the cells are correlated.

Two implementation choices are exposed rather than buried, because both move
results.  First, the BIC used to select a point on the path is $D + m\ln n$,
and the theory does not fix $n$: the $K$ first-stage estimates are the
observations, while the asymptotics justifying Schwarz's rate run in the
number of units.  The default $n = K$ reproduces the group counts reported in
the paper, and \texttt{nbic()} changes it.  Second, the stopping rule compares
$\Delta(A,B)$ against $\lambda|A||B|$, and the pair product grows as groups
absorb one another.  The effective per-merge threshold is therefore not
monotone in the merge order, and some group counts are unreachable by any
single $\lambda$.  On the application in section~\ref{sec:example} the
thresholds run $0.079$, $0.686$, $0.397$, $1.350$, $2.765$, $3.571$: no
$\lambda$ selects five groups.  \texttt{phdid sensitivity, by(m)} walks the
path one merge at a time instead.

\subsection{The Dirichlet process estimator}

The Bayesian model places a Dirichlet process prior with concentration
$\alpha$ on the partition and a normal base measure on the group effects, and
is fitted by collapsed Gibbs sampling, Algorithm~3 of \citet{Neal2000}.  Each
sweep reassigns one cell at a time with the group effect integrated out; the
collapsed marginal likelihood under the exact $\Sigma$ is computed by the
Woodbury identity, equation~(30) of \citet{AroraWagle2026}.

Because the sampler averages over partitions rather than selecting one, its
intervals include uncertainty about the grouping.  In the paper's simulations
that buys coverage of $0.79$--$0.81$ against $0.57$--$0.62$ for the $\ell_0$
estimator when the partition is genuinely uncertain, at essentially no cost
($0.94$ against $0.95$) when the effects are well separated.  The two
estimators are not rivals: the $\ell_0$ estimator is the fixed-variance MAP of
the same model, differing only in the partition prior.

\subsection{Whether a partition is recoverable at all}

Both estimators will partition pure noise if asked to, and their output looks
no different when they do.  In the second application of
\citet{AroraWagle2026} the $\ell_0$ estimator splits 138 event-study effects
into five tidy bands while a randomization test cannot reject a single common
effect.  \texttt{phdid test} is intended to be run first.  It reports a
common-effect $\chi^2$ test on the pooled deviance, an excess-dispersion
heterogeneity share computed against the exact $\Sigma$, and---when placebo
estimates are supplied---the pre/post noise gauge and within-cell
randomization test of section~5.2.3 of the paper.  Each block is tagged with
its provenance, and the command then states a reading in words rather than
leaving the numbers to speak for themselves.

\section{The \texttt{phdid} command}

\subsection{Syntax}

\begin{stsyntax}
phdid setup
\optional{, \underbar{b}(\it{matname}) \underbar{v}(\it{matname})
cells(\it{namelist}) \underbar{w}eights(\it{matname})
\underbar{nu}nits(\it{\#}) nowarn}
\end{stsyntax}

\begin{stsyntax}
phdid test
\optional{, pre(\it{matname}) nsim(\it{\#}) nocenter seed(\it{\#})}
\end{stsyntax}

\begin{stsyntax}
phdid l0
\optional{, \underbar{sel}ect(bic\textbar lambda\textbar m)
\underbar{l}ambda(\it{\#}) m(\it{\#}) nbic(\it{\#})
\underbar{l}evel(\it{\#}) nopath}
\end{stsyntax}

\begin{stsyntax}
phdid bayes
\optional{, alpha(\it{\#}) mu0(\it{\#}) s0sq(\it{\#}) \underbar{iter}s(\it{\#})
\underbar{burn}in(\it{\#}) thin(\it{\#}) chains(\it{\#})
\underbar{marg}inal(exact\textbar diagonal)
init(singletons\textbar pooled\textbar dispersed)
\underbar{l}evel(\it{\#}) seed(\it{\#}) nococlust}
\end{stsyntax}

\begin{stsyntax}
phdid fit, \underbar{part}ition(\it{numlist})
\optional{\underbar{l}evel(\it{\#}) notable}
\end{stsyntax}

Postestimation subcommands are \texttt{aggregate}, \texttt{partition},
\texttt{rhat}, \texttt{sensitivity}, \texttt{covcheck}, \texttt{enumerate} and
\texttt{graph}; see the help files for their options.

\subsection{Options}

A full list is in the online help.  The options that change results rather
than presentation are \texttt{nbic()} and \texttt{select()} for
\texttt{phdid l0}; \texttt{alpha()}, \texttt{mu0()}, \texttt{s0sq()} and
\texttt{marginal()} for \texttt{phdid bayes}; and \texttt{weights()} for
\texttt{phdid setup}, which affects every aggregate but nothing about the
partition posterior.

\subsection{Stored results}

\texttt{phdid l0}, \texttt{phdid bayes} and \texttt{phdid fit} are
\texttt{eclass} and post \texttt{e(b)} and \texttt{e(V)} over the $K$ cell
effects, so \texttt{lincom}, \texttt{nlcom}, \texttt{test} and
\texttt{estimates} work on them without special handling.  After
\texttt{phdid bayes}, \texttt{e(V)} is the posterior covariance across draws
and therefore already includes uncertainty about the partition; after
\texttt{phdid l0} and \texttt{phdid fit} it is the restricted GLS covariance,
which does not.  The posterior draws themselves are left in Mata, which is
what \texttt{phdid aggregate} reads in order to apply each aggregator
draw-by-draw rather than to the posterior mean.

\section{Example}
\label{sec:example}

We reproduce the first application of \citet{AroraWagle2026}, on the
county-level minimum-wage panel of \citet{CallawaySantAnna2021}.  The first
stage is \texttt{csdid}; any heterogeneity-robust estimator would serve.

\begin{stlog}
. csdid lemp, ivar(countyreal) time(year) gvar(first_treat) ///
>     method(dripw) notyet
. phdid setup
. phdid test
\end{stlog}

The test reports $\chi^2(6) = 25.29$ with $p = 0.0003$, a heterogeneity share
of 80\% and a signal-to-noise ratio of 2.24, and reads this as recoverable
heterogeneity.  The $\ell_0$ estimator then selects four groups:

\begin{stlog}
. phdid l0
. phdid aggregate
\end{stlog}

giving an overall effect of $-0.0388$ with a 95\% interval of $[-0.0614,
-0.0162]$, conditional on that partition.  Averaging over partitions instead,

\begin{stlog}
. set seed 20260301
. phdid bayes, alpha(1) iters(20000) burnin(2000) chains(4) init(dispersed)
. phdid rhat
. phdid aggregate
\end{stlog}

gives $-0.0238$ with $[-0.0476, -0.0000]$ and a posterior mean of 2.20 groups.
The two differ because the $\ell_0$ interval conditions on a partition that
carries little posterior mass; the co-clustering matrix, drawn by
\texttt{phdid graph coclust}, shows which pairs of cells the data really do
bind together and which the point partition merely happened to join.

With only seven cells the posterior can be computed exactly:

\begin{stlog}
. phdid enumerate, alpha(1)
\end{stlog}

enumerating all $B_7 = 877$ partitions and returning a posterior mean of
$2.1953$ groups and an overall effect of $-0.0240$ with $[-0.0477, -0.0003]$.
That the sampler reproduced these is the check in section~\ref{sec:valid}.

\section{Validation}
\label{sec:valid}

An R package of the same name, implementing the same estimators, is on CRAN.
Everything the estimators do is deterministic given $(\hat\tau, \Sigma)$
except the Gibbs sampler, and Stata's random-number generator is not R's, so
\texttt{phdid} is certified in two layers.

The deterministic parts---the restricted GLS fit, the agglomeration path, BIC
selection, the specification test, and the exact partition posterior by
enumeration---are certified against the R package to machine precision, using
143 reference values written at 17 significant figures.  The Gibbs sampler is
certified against the exact posterior from enumeration, which both languages
compute deterministically and which is the sampler's own target.  That is the
correct benchmark: reproducing R's draws would be meaningless, while
reproducing R's posterior is not.

The certification script, \texttt{cert/phdid\_cert.do}, also checks invariants
that do not depend on R: relabelling a partition does not change the fit,
aggregation weights are a partition of unity within each term, grouping never
raises a cell's sampling variance, and a coarser partition never fits better.
It runs 165 checks.

\section{Acknowledgements}

The methods are joint work by both authors.  Rohan Wagle maintains the
package.  We thank the authors of \texttt{csdid} and \texttt{drdid}, on which
the examples depend.

\section{References}

\begin{phang}
Arora, P., and R. Wagle. 2026. A Bayesian approach to partial homogeneity in
staggered difference-in-difference. Ashoka University Economics Discussion
Paper 166. Also SSRN Working Paper 7207083,
\url{https://doi.org/10.2139/ssrn.7207083}.
\end{phang}

\begin{phang}
Callaway, B., and P. H. C. Sant'Anna. 2021. Difference-in-differences with
multiple time periods. \emph{Journal of Econometrics} 225: 200--230.
\end{phang}

\begin{phang}
Neal, R. M. 2000. Markov chain sampling methods for Dirichlet process mixture
models. \emph{Journal of Computational and Graphical Statistics} 9: 249--265.
\end{phang}

\begin{phang}
Wade, S., and Z. Ghahramani. 2018. Bayesian cluster analysis: Point estimation
and credible balls. \emph{Bayesian Analysis} 13: 559--626.
\end{phang}

\begin{phang}
Wooldridge, J. M. 2025. Two-way fixed effects, the two-way Mundlak regression,
and difference-in-differences estimators. \emph{Empirical Economics} 69:
2545--2587.
\end{phang}

\aboutauthors

Rohan Wagle is an undergraduate in the Department of Economics at Ashoka
University.  He maintains \texttt{phdid} for both Stata and R.

Parush Arora is in the Department of Economics at Ashoka University.

\end{document}
